\documentclass[11pt]{article}

\usepackage{geometry}
\geometry{a4paper, margin=2.4cm}

\usepackage{fontspec}
\usepackage{amsmath}
\usepackage{amsthm}
\usepackage{unicode-math}
% Polices : Latin Modern (défauts fontspec / unicode-math) — fournis avec tectonic.

\usepackage{polyglossia}
\setdefaultlanguage{french}

\usepackage{microtype}
\usepackage[dvipsnames]{xcolor}
\usepackage[colorlinks=true, linkcolor=RoyalBlue, urlcolor=RoyalBlue]{hyperref}
\usepackage{enumitem}

\setlength{\parindent}{0pt}
\setlength{\parskip}{0.45em}
\setlength{\emergencystretch}{3em}
\setlist{leftmargin=1.6em}

\theoremstyle{plain}
\newtheorem{theoreme}{Théorème}
\newtheorem{proposition}{Proposition}
\renewcommand{\proofname}{Démonstration}

\newcommand{\tp}{^{\top}}
\newcommand{\bc}{\mathbf{c}}
\newcommand{\bb}{\mathbf{b}}
\newcommand{\bG}{\mathbf{G}}
\newcommand{\bZ}{\mathbf{Z}}
\newcommand{\bA}{\mathbf{A}}
\newcommand{\bd}{\mathbf{d}}
\newcommand{\bu}{\mathbf{u}}
\newcommand{\by}{\mathbf{y}}
\newcommand{\bs}{\mathbf{s}}
\newcommand{\bw}{\mathbf{w}}
\newcommand{\bze}{\mathbf{z}}
\newcommand{\bone}{\mathbf{1}}

\title{\bfseries Sur la taille du support de l'optimum OCS génomique}
\author{\textbf{Adel Kaleche}\\[2pt] \small\itshape Affiliation à compléter. \\ \small\itshape Version française ; note de référence en anglais.}
\date{}

\begin{document}
\maketitle

\begin{abstract}
\noindent
Note compagne du solveur support-first. Pour la sélection à contribution optimale (OCS)
génomique, on encadre la taille du support optimal : une borne propre indépendante de $n$
tient dans la limite sans ridge (Théorème~\ref{thm:eps0}), aucune borne universelle ne
survit au ridge (Théorème~\ref{thm:noridge}), et le régime pratique entre les deux est
gouverné par la géométrie conjointe du spectre, des valeurs génétiques et du cap, sans loi
scalaire unique. Une identité KKT (Proposition~\ref{prop:kkt}) relie les deux théorèmes.
Toutes les affirmations empiriques sont reproductibles.
\end{abstract}

\section{La question}

L'OCS génomique résout, pour $n$ candidats de valeurs génétiques $\bb \in \mathbb{R}^n$ et de
matrice de parenté $\bG = \bZ\bZ\tp/s + \varepsilon\mathbf{I}$ ($\bZ$ la matrice $n\times m$ des
génotypes centrés, $s$ une échelle, $\varepsilon$ un petit ridge assurant la définie-positivité),
\begin{equation}\label{eq:ocs}
  \max\ \bb\tp\bc \quad\text{s.c.}\quad \bA\bc = \bd,\ \ \mathbf{0} \le \bc \le \bu,\ \ \bc\tp\bG\bc \le k,
\end{equation}
où $\bA$ est la matrice de budget $q\times n$ ($q=1$ pour le simplexe $\bone\tp\bc=1$ ; $q=2$
pour la forme sexuée $\sum_{\text{mâles}} = \sum_{\text{femelles}} = \tfrac12$). Empiriquement
l'optimum $\bc^\star$ active très peu de candidats : son support $S = \{i : c^\star_i > 0\}$
est $\approx 15$--$30$ sur les panels réels et reste borné quand $n$ croît à cap fixé. Qu'est-ce
qui borne $|S|$ ?

La réponse est un \emph{encadrement}. Dans la limite sans ridge, une borne propre indépendante
de $n$ existe (\S\ref{sec:eps0}) ; dès que le ridge est présent, il n'y a prouvablement aucune
borne universelle, le support pouvant atteindre toute la population (\S\ref{sec:noridge}) ; et
le régime pratique entre les deux est fixé par la géométrie conjointe de $(\text{spectre}, \bb, k)$,
qui n'admet aucune loi scalaire unique et que l'on cartographie empiriquement
(\S\ref{sec:empirical}). Une identité KKT (\S\ref{sec:kkt}) relie les deux théorèmes.

\section{La borne \texorpdfstring{$\varepsilon = 0$}{eps=0}}\label{sec:eps0}

\begin{theoreme}\label{thm:eps0}
Pour $\varepsilon = 0$ (donc $\bG = \bG_0 = \bZ\bZ\tp/s$ de rang $r \le m$),
\eqref{eq:ocs} admet un optimum $\bc^\star$ avec $|S| \le q + r + 1$, indépendant de $n$.
\end{theoreme}

\begin{proof}
Par dualité forte (Slater tient pour toute instance intérieurement réalisable), il existe un
multiplicateur $\lambda^\star \ge 0$ tel que $\bc^\star$ maximise le lagrangien
$L(\bc) = \bb\tp\bc - \lambda^\star\,\bc\tp\bG_0\bc$ sur le polytope
$\{\bA\bc = \bd,\ \bc \ge 0\}$. Comme $\bc\tp\bG_0\bc = \lVert\bZ\tp\bc\rVert^2/s$, la valeur
$L(\bc)$ ne dépend de $\bc$ qu'à travers le couple $(\bZ\tp\bc,\ \bb\tp\bc) \in \mathbb{R}^{r+1}$.
Donc la tranche
\[
  P = \{\, \bc \ge 0 : \bA\bc = \bd,\ \bZ\tp\bc = \bZ\tp\bc^\star,\ \bb\tp\bc = \bb\tp\bc^\star \,\}
\]
est \emph{entièrement optimale} : tout $\bc \in P$ est réalisable ($\bc\tp\bG_0\bc = \lVert\bZ\tp\bc^\star\rVert^2/s \le k$)
et atteint la valeur optimale $\bb\tp\bc^\star$. $P$ est non vide (il contient $\bc^\star$) et
borné (le budget avec $\bc \ge 0$), et il est découpé par $q + r + 1$ lignes d'égalité ; il
possède donc un sommet, et un sommet de $\{\bc \ge 0 : \mathbf{M}\bc = \mathbf{e}\}$ à $q+r+1$
lignes a au plus $q + r + 1$ composantes non nulles. Ce sommet est un optimum de support annoncé.
\end{proof}

C'est l'idée de Carathéodory / Barvinok--Pataki appliquée correctement : on ne compte
\emph{pas} sur la frontière courbe de l'ellipsoïde (où, $\bG$ étant définie positive, tout point
est déjà un point extrême de l'ensemble réalisable, donc le support n'est pas borné) --- on
\emph{fige plutôt l'image de bas rang $\bZ\tp\bc$ et l'objectif}, ce qui rend la tranche optimale
affine, et on y compte les sommets LP. Les solveurs à ensemble actif (support-first ; l'algorithme
de la ligne critique) renvoient un tel sommet, tandis que les solveurs à point intérieur et ADMM
renvoient des points intérieurs non parcimonieux qu'ils seuillent a posteriori --- la borne porte
donc précisément sur ce que calcule un solveur à ensemble actif.

\section{Aucune borne universelle pour \texorpdfstring{$\varepsilon > 0$}{eps>0}}\label{sec:noridge}

\begin{theoreme}\label{thm:noridge}
Pour $\varepsilon > 0$, aucune borne sur $|S|$ de la forme $f(q, r)$ indépendante de $n$
n'existe : le support peut valoir $n$.
\end{theoreme}

\begin{proof}[Démonstration (contre-exemple)]
Prenons $\bG = \varepsilon\mathbf{I}$ --- le cas dégénéré $m = 0$ (aucun marqueur, toute paire
également non apparentée ; $\operatorname{rang}(\bG_0) = 0$). Alors \eqref{eq:ocs}, forme
simplexe, devient
\[
  \max\ \bb\tp\bc \quad\text{s.c.}\quad \bone\tp\bc = 1,\ \bc \ge 0,\ \varepsilon\lVert\bc\rVert^2 \le k,
\]
et $\lVert\bc\rVert^2 = \sum_i c_i^2$ est exactement le proxy du taux de consanguinité. Le minimum
de $\lVert\bc\rVert^2$ sur le simplexe est $1/n$, atteint uniquement au plan uniforme
$\bc = \bone/n$. Pour $k$ légèrement supérieur à $\varepsilon/n$, l'ensemble réalisable est le
simplexe intersecté avec une boule qui se contracte sur $\bone/n$, forçant $|S| = n$.
Numériquement (\texttt{bound\_validation.py}, bloc~7, $n = 300$) : $|S| = 300, 295, 256, 151, 52, 12$
quand $k$ se desserre de $1{,}05\times$ à $50\times$ le minimum. Donc $|S|$ parcourt $[1, n]$ sans
plafond indépendant de $n$.
\end{proof}

Génétiquement, $\bG = \varepsilon\mathbf{I}$ est la limite \emph{sans structure} ; sans schéma
d'apparentement, le cap de diversité ne peut être satisfait qu'en étalant les contributions sur
toute la population. Le petit support observé sur données réelles est donc une propriété de la
\emph{structure} de $\bG$ (et de $\bb$), pas une garantie de pire cas.

\section{Le pont : une identité KKT}\label{sec:kkt}

\begin{proposition}\label{prop:kkt}
À un optimum où la contrainte de parenté est active, les contributions sur le support sont une
fonction affine du feature augmenté $(b_i, \bze_i) \in \mathbb{R}^{m+1}$ de chaque candidat
($\bze_i$ la $i$-ème ligne de $\bZ$) : $\ c^\star_i = \alpha\,b_i + \bw\tp\bze_i + \beta_{\mathrm{sexe}(i)}$
pour $i \in S$.
\end{proposition}

\begin{proof}
Les conditions KKT donnent $\mu \in \mathbb{R}^q$, $\lambda \ge 0$, $\bs \ge 0$ avec
$\bb = \bA\tp\mu + 2\lambda\bG\bc^\star - \bs$ et $s_i c^\star_i = 0$. Sur $S$ ($s_i = 0$) :
$b_i = (\bA\tp\mu)_i + 2\lambda(\bG\bc^\star)_i$. En posant $\by := \bZ\tp\bc^\star$ et en
utilisant $\bG\bc^\star = \bZ\by/s + \varepsilon\bc^\star$, on obtient, pour $i \in S$,
$\varepsilon c^\star_i = (b_i - (\bA\tp\mu)_i)/(2\lambda) - (\bze_i\tp\by)/s$, soit
$c^\star_i = \alpha\,b_i + \bw\tp\bze_i + \beta_{\mathrm{sexe}(i)}$ avec $\alpha = 1/(2\lambda\varepsilon)$,
$\bw = -\by/(s\varepsilon)$, $\beta_{\mathrm{sexe}(i)} = -(\bA\tp\mu)_i/(2\lambda\varepsilon)$.
\end{proof}

Les contributions du support vivent donc sur une famille à $(m+2)$ paramètres. Quand
$\varepsilon \to 0$, le terme $\varepsilon c^\star_i$ s'annule et l'identité force
$\bb - \bA\tp\mu$ dans l'espace-ligne de rang $r$ de $\bZ$, ce qui retrouve le
Théorème~\ref{thm:eps0} et fait apparaître $m$ (le nombre de marqueurs) comme la dimension
pertinente. Pour $\varepsilon > 0$, la même identité ne borne \emph{pas} $|S|$, en accord avec le
Théorème~\ref{thm:noridge}.

\section{Le régime empirique entre les théorèmes}\label{sec:empirical}

Entre les deux énoncés propres se trouve le régime que les sélectionneurs utilisent, cartographié
sur des spectres synthétiques, des alignements de $\bb$ et des caps, ainsi que sur les panels
réels :

\begin{itemize}
\item \textbf{L'alignement de $\bb$ avec le spectre gouverne $|S|$, inversement.} À $\bG$ et $k$
  fixés, placer $\bb$ sur la direction propre dominante (chère en coancestrie) force un grand
  support (moyenne $|S| \approx 146$ sur $n = 800$ pour $\bb$ le vecteur propre de tête) ; étaler
  $\bb$ sur de nombreuses directions bon marché l'effondre à $\approx 4$. Le gain cherché dans une
  direction chère doit être dilué sur beaucoup de candidats pour respecter le cap
  (\texttt{bound\_balign.py}).
\item \textbf{Aucun scalaire unique ne prédit $|S|$ à travers les régimes.} Ni
  $\operatorname{rang}(\bG_0)$, ni le rang effectif / ratio de participation du spectre, ni le coût
  directionnel $\bb\tp\bG\bb/\bb\tp\bb$, ni le support de la direction non contrainte
  $(\bG^{-1}\bb)_+$ ne suit $|S|$ à travers les instances synthétiques et réelles ; la meilleure
  combinaison sans dimension, $|S| \sim (\bb\tp\bG\bb/\bb\tp\bb \cdot 1/k)^{0{,}8}$, n'atteint que
  $R^2 \approx 0{,}58$ et rate le support \emph{relatif} des panels réels (\texttt{bound\_predictor.py},
  \texttt{bound\_lawfit.py}).
\item \textbf{La courbe $|S|(k)$ est une loi de puissance par instance, d'exposant non universel.}
  $|S| \sim k^{-\alpha}$ ajuste bien chaque instance ($R^2 \approx 0{,}9$--$0{,}98$ là où le support
  est bien résolu), avec $\alpha \approx 1$ sur les deux panels réels mais variant de $0{,}5$ à
  $1{,}7$ selon les spectres et alignements synthétiques (\texttt{bound\_curve.py}).
\item \textbf{Panels réels.} Blé ($n = 599$) et souris ($n = 1814$) : $|S| = 25/26$ et $17/26$
  (simplexe / optiSel sexé) au cap de travail, tous deux de l'ordre de la dizaine --- petits et
  stables en $n$, comme le régime à spectre décroissant le prédit qualitativement, sans qu'aucune
  formule n'en fixe la valeur (\texttt{bound\_real.py}).
\end{itemize}

Le Théorème~\ref{thm:noridge} explique \emph{pourquoi} aucune loi scalaire n'existe : il n'y a pas
de borne universelle à prédire, donc la valeur pratique est une fonction de la géométrie complète
$(\text{spectre}, \bb, k)$.

\section{Une borne conditionnelle : l'obstruction et la conjecture}

La question universelle est close négativement (Théorème~\ref{thm:noridge}) ; la question utile est
\emph{conditionnelle}. Deux bornes propres encadrent le reste : à \emph{cap lâche}
($\bc\tp\bG\bc^\star < k$) l'ensemble actif est le seul polytope $\{\bA\bc=\bd,\ \bc\ge0\}$ et
$\bc^\star$ est un sommet LP, donc $|S| \le q$ ; à \emph{rang exact} $d$, le Théorème~\ref{thm:eps0}
donne $|S| \le q + d + 1$.

Une borne en \emph{rang effectif} --- l'objectif réaliste, puisque les vrais spectres décroissent ---
ne découle pas du même argument, et l'obstruction est précise. Le Théorème~\ref{thm:eps0} linéarise la
tranche optimale en épinglant $\bZ\tp\bc$, soit \emph{toute} l'image de $\bG_0$ ; pour un $\bG$ de
plein rang il faudrait épingler les $n$ projections $u_i\tp\bc$, ce qui ne donne que le trivial
$|S| \le q + n + 1$. On ne peut pas n'épingler que les $d$ directions de tête et négliger la queue, car
celle-ci contribue $\bc\tp E\bc = \sum_{i>d} \lambda_i (u_i\tp\bc)^2$ --- une \emph{vraie quadratique},
pas un terme de bas rang --- donc la fixer n'est pas une contrainte linéaire et la tranche n'est plus
un polytope. La courbure de la queue est exactement ce qui laisse le support s'étaler ; le rang
effectif ne peut pas entrer dans un comptage de sommets. Une borne conditionnelle exige donc un autre
argument (perturbation spectrale).

Elle exige aussi deux hypothèses, \emph{toutes deux nécessaires} :
\begin{itemize}
\item \textbf{(A1) Pas de grand plateau dégénéré bon marché.} Un bloc de nombreuses petites valeurs
  propres quasi-égales est un sous-espace bon marché où l'optimum s'étale ; $\bG = \varepsilon\mathbf{I}$
  est le cas extrême et donne $|S| = n$ (Théorème~\ref{thm:noridge}). Un spectre régulièrement
  décroissant n'a pas un tel plateau.
\item \textbf{(A2) $\bb$ évite les directions dominantes.} Même avec un gap, $\bb$ sur le vecteur
  propre de tête (cher) force la dilution et un grand support (moyenne $|S| \approx 146$ sur $800$ dans
  le balayage d'alignement).
\end{itemize}
Les deux contre-exemples sont établis : aucune hypothèse n'est superflue. Sous (A1)+(A2) --- le régime
des vraies matrices génomiques, où $|S| \approx 15$--$30$ --- on conjecture $|S|$ petit et indépendant
de $n$, de l'ordre des directions propres dominantes que $\bb$ excite réellement ; la moitié croissante
($|S|$ augmentant quand $k \to 0$) suit la filiation
contributions\,$\leftrightarrow\,\Delta F\,\leftrightarrow\,N_e$ (Wray \& Thompson 1990 ;
Woolliams \& Bijma 2000). La prouver --- coupler le gap spectral à la projection de $\bb$ sur le
sous-espace dominant, en résistant à la courbure de la queue --- est le problème ouvert, point de
rencontre de l'optimisation (la face du cône perturbé) et de la génétique quantitative (lignées
effectives).

\textbf{Une prise exacte sur le support effectif.} Écrivons le nombre \emph{effectif} de contributeurs
comme le ratio de participation $\mathrm{PR}(\bc^\star) = 1/\lVert\bc^\star\rVert^2$ (égal à $|S|$ pour
des contributions uniformes, moindre si inégales). Au cap actif avec $\bone\tp\bc^\star = 1$, une
identité tient : $\mathrm{PR}(\bc^\star) = R(\bc^\star)/k$, où
$R(\bc^\star) = \bc^{\star\top}\bG\bc^\star / \bc^{\star\top}\bc^\star$ est le quotient de Rayleigh
propre de l'optimum (car $R\lVert\bc^\star\rVert^2 = \bc^{\star\top}\bG\bc^\star = k$). Donc
$\mathrm{PR}(\bc^\star) \le \lambda_1/k$ sans condition, et la borne conditionnelle \emph{se réduit à
borner} $R(\bc^\star)$ --- le coût directionnel de coancestrie que l'optimum est forcé d'encaisser.
C'est là qu'agit (A2) : quand $\bb$ évite les directions dominantes, $\bc^\star$ vit dans la partie bon
marché du spectre et $R(\bc^\star)$ est petit (mesuré : $R \approx 0{,}02$--$0{,}4$ sur les cas
structurés contre $\lambda_1 \approx 45$--$143$, donnant $\mathrm{PR} \approx 2$--$12$), tandis que
$\bG = \varepsilon\mathbf{I}$ force $R = \varepsilon = \lambda_1$ et $\mathrm{PR} = n$. La pièce ouverte
résiduelle est désormais nette et scalaire --- borner $R(\bc^\star)$ a priori sous (A1)+(A2). (La
cardinalité brute $|S|$ vaut $\approx 1{,}5$--$2{,}5\times \mathrm{PR}$ ici ; $\mathrm{PR}$ est la
quantité à loi propre.) Vérifié dans \texttt{bound\_kernel.py}.

\section*{Évidence numérique (reproductible)}

\texttt{bound\_validation.py} (la borne $\varepsilon=0$, la n-indépendance, les balayages ridge et
décroissance spectrale, et le contre-exemple du Théorème~\ref{thm:noridge}) ;
\texttt{bound\_real.py} (blé et souris, après \texttt{research/repro/*\_export.R}) ;
\texttt{bound\_balign.py} (le balayage d'alignement de $\bb$) ; \texttt{bound\_predictor.py} et
\texttt{bound\_lawfit.py} (la recherche de prédicteur / de loi) ; \texttt{bound\_curve.py} (l'ajustement
de $|S|(k)$ par instance). Journal complet : \texttt{research/support\_bound\_sketch.md}.

\section*{Références}
\sloppy
\begin{enumerate}[leftmargin=2em]
\item \textbf{Pataki, G. (1998).} On the rank of extreme matrices in semidefinite programs and the multiplicity of optimal eigenvalues. \emph{Math.\ Oper.\ Res.} 23(2):339--358. DOI 10.1287/moor.23.2.339.
\item \textbf{Markowitz, H. (1956).} The optimization of a quadratic function subject to linear constraints. \emph{Naval Res.\ Logist.\ Q.} 3(1--2):111--133. DOI 10.1002/nav.3800030110.
\item \textbf{Wray, N.R. \& Thompson, R. (1990).} Prediction of rates of inbreeding in selected populations. \emph{Genet.\ Res.} 55(1):41--54. DOI 10.1017/S0016672300025180.
\item \textbf{Woolliams, J.A. \& Bijma, P. (2000).} Predicting rates of inbreeding in populations undergoing selection. \emph{Genetics} 154(4):1851--1864. DOI 10.1093/genetics/154.4.1851.
\item \textbf{Yamashita, M., Mullin, T.J. \& Safarina, S. (2018).} An efficient second-order cone programming approach for optimal selection in tree breeding. \emph{Optim.\ Lett.} 12(7):1683--1697. DOI 10.1007/s11590-018-1229-y.
\item \textbf{Waldmann, P. (2025).} Genomic optimum contribution selection and mate allocation using JuMP. \emph{Bioinform.\ Adv.} 5(1):vbaf259. DOI 10.1093/bioadv/vbaf259.
\end{enumerate}

\end{document}
